% Defense Slides
% DIE-HARL: Durable, Implicit, and Extensible Heterogeneous Agent Reinforcement Learning
% Capt Brandon Hosley, USAF | AFIT-ENS-DS-26-M-094

\documentclass{afitdefense}
\usepackage{booktabs}
\usepackage{multirow}
\usepackage{pifont}   % \ding symbols for checkmarks
\usepackage{nicefrac} % \nicefrac{1}{n} inline fractions
\usetikzlibrary{positioning, shapes.geometric}

\graphicspath{
    {../Chapters/C1/Figures/}
    {../Chapters/C2/Figures/}
    {../Chapters/C3/Figures/}
    {../Figures/}
    {./}
}

% Convenience: colored bold inline text
\newcommand{\term}[1]{\textbf{\textit{#1}}}
\newcommand{\accentbold}[1]{\textcolor{light}{\textbf{#1}}}
\newcommand{\yes}{\textcolor{green}{\ding{51}}}  % checkmark
\newcommand{\done}{\textbf{\textcolor{green}{Complete}}}

\pretitle{Ph.D.\ Defense \textbullet\ AFIT-ENS-DS-26-M-094}
\title{DIE-HARL}
\subtitle{Durable, Implicit, and Extensible\\Heterogeneous Agent Reinforcement Learning}
\author{Capt Brandon Hosley, USAF\\[0.5ex]
    \footnotesize
    Dr.\ Bruce Cox (Chair)\qquad
    Dr.\ Matthew Robbins\qquad
    Maj Nicholas Yielding, Ph.D.}

\begin{document}
\maketitle


%% ==========================================================================
%% OVERVIEW
%% ==========================================================================

\chapter{Overview}

%% ------------------------------------------------------------------
\section[Roadmap]{Dissertation Roadmap}

Three questions, one unified answer.

\begin{center}
\renewcommand{\arraystretch}{1.6}
\begin{tabular}{@{}C{0.3cm}L{3.5cm}L{4.5cm}L{4.0cm}@{}}
    \cellcolor{dark}\textcolor{light}{\textbf{\#}} &
    \cellcolor{dark}\textcolor{white}{\textbf{Contribution}} &
    \cellcolor{dark}\textcolor{white}{\textbf{Central Question}} &
    \cellcolor{dark}\textcolor{white}{\textbf{Thread}} \\[0.3em]
    \textbf{C1} & \textbf{Policy Upsampling} &
        Can pretraining \textit{smaller} teams reduce training cost? &
        \textit{Efficiency} \\
    \textbf{C2} & \textbf{Implicit Indication} &
        Can one policy serve \textit{heterogeneous} agents without IDs? &
        \textit{Flexibility} \\
    \textbf{C3} & \textbf{Paradigm Comparison} &
        Representational or architectural invariance---which wins? &
        \textit{Validation} \\
\end{tabular}
\end{center}

\vspace{0.5ex}
\shadedbox{Unifying thesis: HARL efficiency and flexibility are fundamentally
\textbf{representational} problems before they are architectural ones.}

%% ------------------------------------------------------------------
\section[Contributions Status]{Contributions at a Glance}

\begin{center}
\renewcommand{\arraystretch}{1.5}
\begin{tabular}{@{}L{2.5cm}L{4.5cm}L{3.5cm}C{2.0cm}@{}}
    \cellcolor{dark}\textcolor{white}{\textbf{Contribution}} &
    \cellcolor{dark}\textcolor{white}{\textbf{Key Finding}} &
    \cellcolor{dark}\textcolor{white}{\textbf{Environments}} &
    \cellcolor{dark}\textcolor{white}{\textbf{Status}} \\[0.3em]
    \textbf{C1:} Policy Upsampling &
        Environment-dependent benefit; \textit{task structure} governs curriculum utility &
        Waterworld, Multiwalker, LBF &
        \done \\
    \textbf{C2:} Implicit Indication &
        Matches HAPPO; $\nicefrac{1}{|I|}$ storage; \textit{zero-shot} robustness &
        HyperGrid &
        \done \\
    \textbf{C3:} Paradigm Comparison &
        Masking outperforms GNNs by \textbf{2.3--4.4$\times$} &
        HyperGrid &
        \done \\
\end{tabular}
\end{center}

\vspace{0.5ex}
\textit{All three contributions complete and integrated into a single coherent framework.}


%% ==========================================================================
%% MOTIVATION
%% ==========================================================================

\chapter{Motivation}

%% ------------------------------------------------------------------
\section{The Autonomous Systems Imperative}\twocolumn\raggedright

\subsection{Deployments We Can Imagine}

\begin{itemize}
    \item 2023: Deputy Sec.\ Hicks announces the \textbf{Replicator Initiative}---hundreds of autonomous systems fielded in \textit{months, not years}
    \item DARPA \textbf{OFFSET}: 100+ agent swarm demonstrations
    \item Robotic warehouses, mixed-sensor drone fleets, agricultural robotics---heterogeneous multi-agent systems are \textit{already here}
\end{itemize}

\subsection{The Gap}

\begin{itemize}
    \item Capability exists in demonstration
    \item \textbf{Training cost} remains the primary bottleneck to operational deployment
    \item Heterogeneous fleets break the assumptions of standard MARL methods
\end{itemize}

\pagebreak

\begin{Center}
\vspace{1em}
\includegraphics[width=0.9\columnwidth, keepaspectratio]{drone-mid.pdf}
\vspace{0.5em}

\scriptsize\textit{Heterogeneous autonomous systems require more than\\
homogeneous assumptions can provide.}
\end{Center}

%% ------------------------------------------------------------------
\section{Heterogeneity: The Operational Reality}\twocolumn\raggedright

\subsection{Two Kinds of Heterogeneity}

\begin{itemize}
    \item \term{Behavioral heterogeneity}: same structure, different \textit{learned} roles
    \begin{itemize}
        \item Standard parameter sharing handles this well
    \end{itemize}
    \item \term{Intrinsic (structural) heterogeneity}: different sensors, effectors, or observation types
    \begin{itemize}
        \item Shared policies break down---observation spaces \textit{don't align}
    \end{itemize}
\end{itemize}

\subsection{The Status Quo}

\begin{itemize}
    \item Explicit agent IDs (one-hot vectors)
    \item Per-type separate policies
    \item Architectural solutions (attention, GNNs)
    \item Each carries costs in storage, flexibility, or sample efficiency
\end{itemize}

\pagebreak

\begin{Center}
\renewcommand{\arraystretch}{1.4}
\begin{tabular}{@{}L{2.5cm}L{3.0cm}@{}}
    \cellcolor{dark}\textcolor{white}{\textbf{Approach}} &
    \cellcolor{dark}\textcolor{white}{\textbf{Cost}} \\
    One-hot IDs & Brittle at deployment \\
    Per-type policies & $|I|\times$ storage \\
    Attention / GNNs & Architectural complexity \\
    \textbf{Implicit Indication} & \textcolor{green}{\textbf{None of the above}} \\
\end{tabular}
\end{Center}

%% ------------------------------------------------------------------
\section{Three Research Challenges}

\begin{center}
\renewcommand{\arraystretch}{1.8}
\begin{tabular}{@{}C{0.4cm}L{4.0cm}L{7.5cm}@{}}
    \cellcolor{dark}\textcolor{light}{\textbf{1}} &
    \cellcolor{dark}\textcolor{white}{\textbf{Efficiency}} &
    \cellcolor{dark}\textcolor{white}{Training large heterogeneous teams is expensive.
        Can \textit{curricula} help?} \\
    \textbf{C1} &
        Policy Upsampling &
        Pretrain smaller teams; scale up via \textbf{policy duplication} \\[0.3em]
    \cellcolor{dark}\textcolor{light}{\textbf{2}} &
    \cellcolor{dark}\textcolor{white}{\textbf{Flexibility}} &
    \cellcolor{dark}\textcolor{white}{Structural heterogeneity breaks parameter sharing.
        Can \textit{representation} fix it?} \\
    \textbf{C2} &
        Implicit Indication &
        Shared policy via \textbf{observation homogenization} \\[0.3em]
    \cellcolor{dark}\textcolor{light}{\textbf{3}} &
    \cellcolor{dark}\textcolor{white}{\textbf{Validation}} &
    \cellcolor{dark}\textcolor{white}{When we solve heterogeneity,
        which paradigm actually wins?} \\
    \textbf{C3} &
        Paradigm Comparison &
        Representational vs.\ \textbf{architectural} invariance \\
\end{tabular}
\end{center}

%% ------------------------------------------------------------------
\section{Unifying Thesis}

\begin{Center}
    \large
    Training efficiency, heterogeneity handling, and deployment
    flexibility in HARL are fundamentally\\[1ex]
    \LARGE\textbf{representational} problems\\[1ex]
    \large before they are \textbf{architectural} ones.

    \vspace{1.5ex}
    \normalsize
    \textit{Task structure and heterogeneity type determine optimal method selection;\\
    good representational design yields robustness as a natural consequence.}
\end{Center}


%% ==========================================================================
%% TECHNICAL FOUNDATIONS
%% ==========================================================================

\chapter{Technical Foundations}

%% ------------------------------------------------------------------
\section{Multi-Agent RL Foundations}\twocolumn\raggedright

\subsection{From Single-Agent to Multi-Agent}

\begin{itemize}
    \item \term{MDP}: states $s$, actions $a$, rewards $r$, transitions $\mathcal{T}$,
        policy $\pi$
    \item \term{Dec-POMDP}: each agent $i$ observes only $o_i \subset s$;
        acts on local information alone
    \item Core challenge: \term{non-stationarity}---the environment appears
        non-stationary to each agent as others learn simultaneously
\end{itemize}

\subsection{The CTDE Paradigm}

\begin{itemize}
    \item \term{Centralized Training, Decentralized Execution (CTDE)}: train with
        global information; deploy with local policies only
    \item Resolves non-stationarity during training
    \item Preserves real-world deployability
\end{itemize}

\pagebreak

\begin{Center}
\vspace{0.3em}
\begin{tikzpicture}[
    box/.style={draw=dark, fill=dark!10, rounded corners=3pt,
                minimum width=2.4cm, minimum height=0.7cm,
                align=center, font=\scriptsize},
    crit/.style={draw=dark, fill=light!40, rounded corners=3pt,
                minimum width=2.4cm, minimum height=0.7cm,
                align=center, font=\scriptsize\bfseries},
    arr/.style={-stealth, thick, dark}]

    % Critic (centralized)
    \node[crit] (critic) at (0, 0) {Centralized\\Critic $V(s)$};
    % Actors
    \node[box] (a1) at (-2.2, -1.5) {Actor $\pi_1(o_1)$};
    \node[box] (a2) at (0,    -1.5) {Actor $\pi_2(o_2)$};
    \node[box] (a3) at ( 2.2, -1.5) {Actor $\pi_i(o_i)$};
    % Global state node
    \node[font=\scriptsize, dark] (gs) at (0, 1.1) {Global state $s$ (training only)};
    % Arrows
    \draw[arr] (gs) -- (critic);
    \draw[arr] (critic) -- (a1);
    \draw[arr] (critic) -- (a2);
    \draw[arr] (critic) -- (a3);
    % Labels
    \node[font=\tiny, dark] at (-2.2, -2.15) {deploy: $o_1$ only};
    \node[font=\tiny, dark] at (0,    -2.15) {deploy: $o_2$ only};
    \node[font=\tiny, dark] at ( 2.2, -2.15) {deploy: $o_i$ only};
\end{tikzpicture}
\vspace{0.2em}

\scriptsize\textit{CTDE: centralized critic (global state, training only)\\
    drives decentralized actors (local obs, deployment)}
\end{Center}

%% ------------------------------------------------------------------
\section{The PPO Family in MARL}\twocolumn\raggedright

\subsection{Algorithm Lineage}

\begin{itemize}
    \item \textbf{PPO} (Schulman 2017): stable on-policy gradient updates
        via clipped surrogate objective
    \item \textbf{MAPPO}: PPO with shared centralized critic across agents
    \item \textbf{IPPO}: independent PPO---each agent trains its own critic
    \item \textbf{HAPPO} (Zhong 2024): \term{Heterogeneous-Agent PPO}---
        separate actor \textit{and} critic per agent type
\end{itemize}

\subsection{Baselines in This Work}

\begin{itemize}
    \item \textbf{HAPPO} serves as the primary comparison in C2 and C3
    \item \textbf{PIC} (Liu 2020): \term{Permutation Invariant Critic}
        using graph neural networks---the architectural baseline in C3
\end{itemize}

\pagebreak

\renewcommand{\arraystretch}{1.5}
\begin{tabular}{@{}L{2.2cm}ccc@{}}
    \cellcolor{dark}\textcolor{white}{\textbf{Algorithm}} &
    \cellcolor{dark}\textcolor{white}{\textbf{Shared $\pi$?}} &
    \cellcolor{dark}\textcolor{white}{\textbf{Shared $V$?}} &
    \cellcolor{dark}\textcolor{white}{\textbf{Heterogeneous?}} \\
    MAPPO  & \yes & \yes & Limited \\
    IPPO   & ---  & ---  & Implicit \\
    HAPPO  & ---  & ---  & \yes \\
    \textbf{Impl.\ Ind.} & \textbf{\yes} & \textbf{\yes} & \textbf{\yes} \\
\end{tabular}

%% ------------------------------------------------------------------
\section{The Heterogeneity Challenge}

\begin{itemize}
    \item \term{Parameter sharing} is the most efficient strategy when it
        applies---one network, all agents
    \item With \textbf{structural heterogeneity}, observation dimensions differ:
        agent $i$ observes $o_i \in \mathbb{R}^{d_i}$, agent $j$ observes
        $o_j \in \mathbb{R}^{d_j}$, $d_i \neq d_j$
    \item Feeding $o_i$ and $o_j$ to the same network is not directly possible
    \item Current solutions force a choice:
    \begin{itemize}
        \item \textit{Explicit IDs} --- one-hot vectors reveal type but are
            brittle under sensor change or team-size variation
        \item \textit{Per-type policies} --- full flexibility but $|I|\times$ storage
            and no generalization across types
        \item \textit{Architectural invariance (GNNs)} --- handles variable teams
            but adds complexity and may not generalize
    \end{itemize}
\end{itemize}

\shadedbox{None of the existing approaches simultaneously achieve
\textbf{shared parameters}, \textbf{structural heterogeneity},
and \textbf{deployment robustness}.}

%% ------------------------------------------------------------------
\section{The Observation-Space Coupling Problem}

\begin{itemize}
    \item In many standard MARL benchmarks, \textbf{observation dimension is
        coupled to team size}: $d_i \propto N$ agents
    \item Adding or removing an agent changes the \textit{input shape} of every
        policy network
    \item This makes policy reuse across team sizes \textbf{architecturally impossible}
        without modification
    \item C1 confronts this directly: policy \textit{upsampling} requires
        \textbf{decoupled observation schemas}
    \begin{itemize}
        \item \term{Truncated schema}: fixed-width observation regardless of team size;
            missing ally slots are zero-padded
        \item \term{Ally-ignorant schema}: agents observe environment state only,
            no ally information at all
    \end{itemize}
\end{itemize}

\shadedbox{Solving team-scaling is not primarily an algorithmic challenge---
it is a \textbf{representational} one.}

%% ------------------------------------------------------------------
\section{Environments Overview}

\begin{center}
\begin{tabular}{C{3.0cm}C{3.0cm}C{3.0cm}C{3.5cm}}
    \cellcolor{dark}\textcolor{white}{\textbf{Waterworld}} &
    \cellcolor{dark}\textcolor{white}{\textbf{Multiwalker}} &
    \cellcolor{dark}\textcolor{white}{\textbf{LBF}} &
    \cellcolor{dark}\textcolor{white}{\textbf{HyperGrid}} \\[0.3em]
    \includegraphics[width=2.8cm, keepaspectratio]{waterworld.png} &
    \includegraphics[width=2.8cm, keepaspectratio]{multiwalker.png} &
    \includegraphics[width=2.8cm, keepaspectratio]{lbf.png} &
    \includegraphics[width=3.2cm, keepaspectratio]{marl-minigrid.png} \\[0.3em]
    Continuous control;\newline homogeneous;\newline low coordination &
    Tight temporal coordination;\newline sparse reward;\newline spatial roles &
    Discrete;\newline dynamic skill heterogeneity;\newline role variation &
    $n$-d MiniGrid;\newline configurable sensors;\newline intrinsic heterogeneity \\[0.3em]
    \multicolumn{3}{c}{\textit{Contribution 1}} & \textit{Contributions 2 \& 3} \\
\end{tabular}
\end{center}


%% ==========================================================================
%% CONTRIBUTION 1
%% ==========================================================================

\chapter[C1]{Contribution 1: Policy Upsampling}

%% ------------------------------------------------------------------
\section[C1: Central Question]{C1---Can We Train Fewer, Then Scale?}\twocolumn\raggedright

\subsection{The Core Hypothesis}

\begin{itemize}
    \item Training an $N$-agent team from scratch is expensive
    \item \textbf{What if we pretrain a smaller $k$-agent team, then duplicate
        policies to fill an $N$-agent team?}
    \item The duplicated policy provides a warm start;
        \term{policy upsampling} completes the team
\end{itemize}

\subsection{Challenges to Address}

\begin{itemize}
    \item Observation-space coupling: adding agents changes input shapes
    \item Requires a \textbf{fixed-width observation schema} (truncated or ally-ignorant)
    \item Does a ``warm-start'' policy \textit{actually} accelerate training---
        or does it hinder adaptation to the larger team?
\end{itemize}

\pagebreak

\vspace{1em}
\includegraphics[width=\columnwidth, height=4.5cm, keepaspectratio]{training_1.png}
\vspace{0.4em}

\scriptsize\textbf{Phase 1:} Pretrain $k$-agent policy.
\textbf{Phase 2:} Duplicate to $N$ agents; continue training.

%% ------------------------------------------------------------------
\section{The Agent-Steps Metric}

\begin{itemize}
    \item Standard metrics (training steps, wall time) \textbf{obscure true cost}
        when team sizes differ
    \item An 8-agent team does 8$\times$ as much work per environment step as a
        single agent
    \item \term{Agent-steps} $= \text{agents} \times \text{training steps}$
        provides a \textbf{normalized cost measure} across configurations
\end{itemize}

\vspace{0.8em}
\begin{center}
\renewcommand{\arraystretch}{1.4}
\begin{tabular}{@{}L{4.5cm}cc@{}}
    \cellcolor{dark}\textcolor{white}{\textbf{Scenario}} &
    \cellcolor{dark}\textcolor{white}{\textbf{Env.\ Steps}} &
    \cellcolor{dark}\textcolor{white}{\textbf{Agent-Steps}} \\
    Train 8 agents from scratch & $T$ & $8T$ \\
    Pretrain 3, fine-tune to 8 &
        $T_{\text{pre}} + T_{\text{ft}}$ &
        $3T_{\text{pre}} + 8T_{\text{ft}}$ \\
    \textbf{Curriculum benefit} & \textbf{same $T$} &
        \textbf{$< 8T$ if $T_{\text{ft}} \ll T$} \\
\end{tabular}
\end{center}

\vspace{0.6em}
\shadedbox{AUC (area under the reward curve) normalized by agent-steps
serves as the primary \textbf{efficiency metric}.}

%% ------------------------------------------------------------------
\section{Observation Schema Design}\twocolumn\raggedright

\subsection{The Problem}

Standard benchmarks encode ally observations in a \textbf{variable-length}
block---adding one ally changes every agent's observation dimension.

\subsection{Two Proposed Schemas}

\begin{itemize}
    \item \term{Truncated schema}: observation has a fixed maximum width;
        present allies fill slots from the front; absent slots are \textbf{zero-padded}
    \begin{itemize}
        \item Preserves ally information up to a maximum team size
    \end{itemize}
    \item \term{Ally-ignorant schema}: ally observations are removed entirely;
        agent observes only its own sensors and task features
    \begin{itemize}
        \item Fully decoupled; limits coordination but eliminates coupling
    \end{itemize}
\end{itemize}

\textit{Both schemas enable the same network to operate at any team size
without architectural changes.}

\pagebreak

\vspace{1em}
\includegraphics[width=\columnwidth, height=4.5cm, keepaspectratio]{training_2.png}

\scriptsize\textit{Fixed-width observation schema: ally slots are filled
or zero-padded, keeping input dimension constant across team sizes.}

%% ------------------------------------------------------------------
\section{C1 Results: Waterworld}\twocolumn\raggedright

\subsection{Environment}

Continuous control; agents collect food and avoid poison pellets;
low coordination requirements; \textbf{homogeneous} agent structure.

\subsection{Findings}

\begin{itemize}
    \item Pretraining \textbf{accelerates convergence} at all target team sizes
    \item Benefit scales with the \textit{ratio} of pretrain size to target size
    \item Largest gains seen for greatest size jump (e.g., $3 \to 8$)
    \item Final performance matches or exceeds baseline
\end{itemize}

\shadedbox{\textbf{Clear curriculum win.} This is the ideal case:
low coordination means the smaller-team policy transfers directly.}

\pagebreak

\includegraphics[width=\columnwidth, height=5cm, keepaspectratio]{Waterworld-AUCs.png}

%% ------------------------------------------------------------------
\section{C1 Results: Multiwalker}\twocolumn\raggedright

\subsection{Environment}

Tight \textbf{temporal coordination}; agents must carry a package cooperatively;
failure of any agent collapses the whole team; sparse reward.

\subsection{Findings}

\begin{itemize}
    \item Without pretraining, large teams ($\geq 6$ walkers) \textbf{frequently fail
        to learn} at all
    \item Pretraining acts as a \term{stabilization mechanism}---it gives the
        team a coordinated foundation to build from
    \item Final performance is similar, but \textit{the path matters}:
        curriculum avoids catastrophic failure modes
\end{itemize}

\shadedbox{\textbf{Curriculum as stability, not just speed.}
Task structure determines the nature of the curriculum benefit.}

\pagebreak

\includegraphics[width=\columnwidth, height=5cm, keepaspectratio]{Multiwalker-AUCs.png}

%% ------------------------------------------------------------------
\section{C1 Results: Level-Based Foraging}\twocolumn\raggedright

\subsection{Environment}

Discrete gridworld; agents have \textbf{skill levels} that determine which items
they can collect; roles emerge dynamically from skill distribution;
\term{intrinsic heterogeneity} via level variation.

\subsection{Findings}

\begin{itemize}
    \item Results are \textbf{limited and inconsistent} across team configurations
    \item Skill-level heterogeneity creates complex role relationships
        that do not transfer cleanly from smaller teams
    \item The ``correct'' team composition at small scale may be suboptimal
        at large scale due to shifting skill distributions
\end{itemize}

\shadedbox{\textbf{Task structure moderates curriculum benefit.}
When heterogeneity \textit{itself} changes with team size,
simple policy duplication cannot bridge the gap.}

\pagebreak

\includegraphics[width=\columnwidth, height=5cm, keepaspectratio]{LBF-AUCs.png}

%% ------------------------------------------------------------------
\section[C1 Synthesis]{C1 Synthesis: What We Learned}

\begin{center}
\renewcommand{\arraystretch}{1.5}
\begin{tabular}{@{}L{2.5cm}L{3.5cm}L{5.5cm}@{}}
    \cellcolor{dark}\textcolor{white}{\textbf{Environment}} &
    \cellcolor{dark}\textcolor{white}{\textbf{Outcome}} &
    \cellcolor{dark}\textcolor{white}{\textbf{Driver}} \\
    Waterworld & Strong acceleration & Low coordination; clean transfer \\
    Multiwalker & Stabilization benefit & Prevents failure; tight coordination \\
    LBF & Limited, inconsistent & Dynamic heterogeneity resists scaling \\
\end{tabular}
\end{center}

\vspace{0.6em}
\textbf{Cross-cutting insight:} The bottleneck in every case is
\textit{representational}---either the observation schema prevents transfer,
or the heterogeneity structure changes with team size.

\vspace{0.6em}
\shadedbox{If we could design observation spaces to \textbf{naturally support}
structural heterogeneity, we would eliminate these barriers at the source.
\textit{That is the motivation for C2.}}


%% ==========================================================================
%% CONTRIBUTION 2
%% ==========================================================================

\chapter[C2]{Contribution 2: Implicit Indication}

%% ------------------------------------------------------------------
\section[C2: Central Question]{C2---Can One Policy Serve Many Agent Types?}\twocolumn\raggedright

\subsection{The Challenge}

\begin{itemize}
    \item Structural heterogeneity: agent $i$ has sensors $\mathcal{C}_i$,
        agent $j$ has sensors $\mathcal{C}_j$, and $\mathcal{C}_i \neq \mathcal{C}_j$
    \item Existing approaches: \textit{explicit} type indicators
        (one-hot vectors, embeddings) fed alongside observations
    \item These are \textbf{fragile}: change the team composition, retrain from scratch
\end{itemize}

\subsection{The Proposal}

\begin{itemize}
    \item \textbf{No explicit type indicators}
    \item Instead, design an observation space that \textit{implicitly} encodes
        agent identity through the pattern of \textbf{which elements are populated}
    \item A single shared policy---one network, all agent types
\end{itemize}

\pagebreak

\begin{Center}
    \vspace{1em}
    \includegraphics[width=0.95\columnwidth, keepaspectratio]{screen_pursuit_channels.png}
    \vspace{0.5em}

    \scriptsize\textit{Agents with different sensor channels occupy
    different regions of the homogenized observation space.}
\end{Center}

%% ------------------------------------------------------------------
\section{The Implicit Indication Framework}

\begin{itemize}
    \item \term{Homogenized observation space $\tilde{\mathcal{O}}$}:
        the union of all agent-specific observation elements
        $\tilde{\mathcal{O}} = \bigcup_{i \in I} \mathcal{O}_i$
    \item Each agent $i$ populates only its relevant elements;
        all other elements are \textbf{masked to zero}
    \item The policy $\pi_\theta(\cdot \mid \tilde{o})$ receives the full
        $\tilde{o}$ vector but infers agent identity from the
        \term{pattern of populated vs.\ masked elements}
    \item \textbf{No explicit type label is needed}---the structure of the
        observation \textit{is} the indicator
    \item \term{Semantic decomposability} (required): each observation element
        must have a \textbf{consistent meaning} across all agent types
    \begin{itemize}
        \item A ``forward-facing camera channel'' means the same thing
            for any agent that has it
        \item A ``thermal sensor channel'' may be absent in an agent with only RGB
    \end{itemize}
\end{itemize}

%% ------------------------------------------------------------------
\section{Semantic Decomposability}\twocolumn\raggedright

\subsection{Formal Requirement}

\begin{itemize}
    \item Observation elements must have \textbf{type-consistent semantics}:
        $\mathcal{O}_i \cap \mathcal{O}_j$ contains elements with the same
        meaning for both agents $i$ and $j$
    \item Elements in $\mathcal{O}_i \setminus \mathcal{O}_j$ are
        simply absent (zeroed) for agent $j$
\end{itemize}

\subsection{Theoretical Guarantee}

\begin{itemize}
    \item Based on \textbf{Deep Sets} theory (Zaheer 2017):
        any collection of functions over heterogeneous input spaces can be
        represented by a \textit{single} function over the homogenized domain
    \item Masking preserves representability: the injective embedding from
        each $\mathcal{O}_i$ into $\tilde{\mathcal{O}}$ is disjoint for
        type-distinct elements
\end{itemize}

\pagebreak

\vspace{0.5em}
\begin{center}
\renewcommand{\arraystretch}{1.5}
\begin{tabular}{@{}L{1.5cm}L{2.0cm}L{2.0cm}@{}}
    \cellcolor{dark}\textcolor{white}{\textbf{Agent}} &
    \cellcolor{dark}\textcolor{white}{\textbf{RGB cam}} &
    \cellcolor{dark}\textcolor{white}{\textbf{Thermal}} \\
    Type A & \yes & --- \\
    Type B & --- & \yes \\
    Type C & \yes & \yes \\
\end{tabular}
\end{center}

\vspace{0.8em}
\shadedbox{The policy \textit{sees} the pattern and infers agent type
without being told explicitly.}

%% ------------------------------------------------------------------
\section[HyperGrid]{HyperGrid: A Purpose-Built Evaluation Environment}\twocolumn\raggedright

\subsection{Why a New Environment?}

\begin{itemize}
    \item Existing benchmarks do not support \textbf{configurable sensor
        heterogeneity} at the level needed for systematic evaluation
    \item A controlled environment is needed to isolate the effect of
        \textit{observation structure} from other confounds
\end{itemize}

\subsection{Design}

\begin{itemize}
    \item MiniGrid-based $n$-dimensional discrete gridworld
    \item Heterogeneity arises from \textbf{sensor channel configuration}:
        which channels an agent observes
    \item Four training configurations tested:
    \begin{itemize}
        \item \term{Complete}: all agents see all channels
        \item \term{Intersecting-span}: overlapping channel subsets
        \item \term{Disjoint-span}: non-overlapping channel subsets
        \item \term{Incomplete}: some channels missing from all agents
    \end{itemize}
\end{itemize}

\pagebreak

\vspace{0.5em}
\includegraphics[width=\columnwidth, height=4.5cm, keepaspectratio]{marl-minigrid.png}

\scriptsize\textit{HyperGrid environment: color-coded sensor channels
determine each agent's observation space.}

%% ------------------------------------------------------------------
\section{C2 Training Results}\twocolumn\raggedright

\subsection{Comparison: Implicit Indication vs.\ HAPPO}

\begin{itemize}
    \item \textbf{HAPPO} (Zhong 2024) maintains separate policies per agent type;
        serves as the strong heterogeneous baseline
    \item Implicit Indication uses a \textbf{single shared policy} across all types
\end{itemize}

\subsection{Finding}

\begin{itemize}
    \item Implicit Indication \textbf{matches HAPPO training performance}
        across all four sensor configurations
    \item Same reward achieved with \textbf{$\frac{1}{|I|}$ storage}
        (one policy instead of $|I|$ per-type policies)
    \item Disjoint-span configuration shows the \textit{strongest} relative
        improvement for implicit indication
\end{itemize}

\pagebreak

\includegraphics[width=\columnwidth, height=5cm, keepaspectratio]{training_curves.png}

%% ------------------------------------------------------------------
\section{C2 Robustness Results}\twocolumn\raggedright

\subsection{Eight Evaluation Conditions}

After training, agents are evaluated under conditions they were
\textit{not trained for}:

\begin{enumerate}
    \item Sensor loss (channel removed)
    \item Sensor degradation (noise added)
    \item Sensor improvement (noise removed)
    \item Coverage reduction (fewer channels)
    \item Coverage expansion (more channels)
    \item Shuffled agent assignments
    \item Incomplete coverage generalization
    \item \textbf{Zero-shot}: novel agent type, unseen during training
\end{enumerate}

\pagebreak

\subsection{Finding}

\includegraphics[width=\columnwidth, height=3.5cm, keepaspectratio]{performance_delta.png}

\vspace{0.5em}
\shadedbox{\textbf{Robustness emerges naturally.}
Implicit Indication handles novel configurations without explicit
robustness training---a \textit{free} benefit of representational design.}

%% ------------------------------------------------------------------
\section[C2 Synthesis]{C2 Synthesis: One Policy, Many Agents}

\begin{center}
\renewcommand{\arraystretch}{1.6}
\begin{tabular}{@{}L{4.5cm}L{7.0cm}@{}}
    \cellcolor{dark}\textcolor{white}{\textbf{Property}} &
    \cellcolor{dark}\textcolor{white}{\textbf{Result}} \\
    Performance vs.\ HAPPO & \textbf{Matches} HAPPO across all configurations \\
    Storage footprint & $\frac{1}{|I|}$ vs.\ per-type policies \\
    Zero-shot generalization & \textbf{Yes}---novel agent types handled gracefully \\
    Sensor change tolerance & \textbf{Yes}---no retraining required \\
    Team-size flexibility & \textbf{Yes}---same policy, variable team size \\
    Requirement & \textit{Semantic decomposability} of observation elements \\
\end{tabular}
\end{center}

\vspace{0.6em}
\shadedbox{C2 demonstrates that the right \textbf{representational choice}
eliminates the need for explicit type indicators entirely.
But is this approach better than architectural alternatives?
\textit{That is the question for C3.}}


%% ==========================================================================
%% CONTRIBUTION 3
%% ==========================================================================

\chapter[C3]{Contribution 3: Paradigm Comparison}

%% ------------------------------------------------------------------
\section[C3: Central Question]{C3---Representation or Architecture?}\twocolumn\raggedright

\subsection{Two Paths to Invariance}

There are two principled ways to handle structural heterogeneity:

\begin{itemize}
    \item \term{Representational invariance}: design the \textit{input space}
        to accommodate all agent types (Implicit Indication---C2)
    \item \term{Architectural invariance}: design the \textit{network architecture}
        to be invariant to agent type and team composition (GNNs, attention)
\end{itemize}

\subsection{The Hypothesis Going In}

\begin{itemize}
    \item GNNs are expressive, flexible, and have strong theoretical properties
    \item PIC (Liu 2020): \term{Permutation Invariant Critic}---
        a purpose-built GNN for heterogeneous MARL
    \item \textit{Prior expectation}: the more sophisticated approach should win
\end{itemize}

\pagebreak

\subsection{The Surprise}

\vspace{1em}
\begin{Center}
\Large\textbf{It doesn't.}

\vspace{1em}
\normalsize
Simple observation masking outperforms\\
graph neural networks \textbf{consistently} and\\
\textbf{substantially} across all conditions.
\end{Center}

%% ------------------------------------------------------------------
\section{Architectural vs.\ Representational Invariance}

\begin{center}
\renewcommand{\arraystretch}{1.6}
\begin{tabular}{@{}L{3.0cm}L{4.5cm}L{4.5cm}@{}}
    \cellcolor{dark}\textcolor{white}{\textbf{Property}} &
    \cellcolor{dark}\textcolor{white}{\textbf{Representational (Impl.\ Ind.)}} &
    \cellcolor{dark}\textcolor{white}{\textbf{Architectural (PIC / GNN)}} \\
    Mechanism & Homogenized obs.\ + masking &
        Graph conv.\ layers + symmetric pooling \\
    Agent info & Inferred from obs.\ pattern &
        Explicit type attributes on GNN nodes \\
    Network type & Standard feed-forward MLP & Graph Neural Network \\
    Variable teams & Via fixed-width obs.\ schema & Via GNN node addition \\
    Complexity & \textbf{Lower} & Higher \\
    Trained with & MAPPO & MAPPO \\
\end{tabular}
\end{center}

\vspace{0.5em}
\textit{Both approaches are given matched network capacities and hyperparameters;
the only difference is the invariance mechanism.}

%% ------------------------------------------------------------------
\section{C3 Experimental Design}

\begin{itemize}
    \item Same \textbf{HyperGrid} environment as C2---controlled, reproducible
    \item \textbf{Four sensor configurations}: complete, intersecting-span,
        disjoint-span, incomplete
    \item \textbf{Eight evaluation conditions} (same as C2): sensor loss/gain,
        coverage changes, shuffled assignments, zero-shot generalization
    \item Matched network capacities; identical hyperparameter sweeps
\end{itemize}

\vspace{0.5em}
\begin{center}
\renewcommand{\arraystretch}{1.4}
\begin{tabular}{@{}L{3.5cm}cc@{}}
    \cellcolor{dark}\textcolor{white}{\textbf{Controlled Variable}} &
    \cellcolor{dark}\textcolor{white}{\textbf{Impl.\ Ind.}} &
    \cellcolor{dark}\textcolor{white}{\textbf{PIC}} \\
    Network capacity & Matched & Matched \\
    Training algorithm & MAPPO & MAPPO \\
    Hyperparameters & Identical sweep & Identical sweep \\
    Environment & HyperGrid & HyperGrid \\
    \textbf{Invariance mechanism} & \textbf{Representational} & \textbf{Architectural} \\
\end{tabular}
\end{center}

%% ------------------------------------------------------------------
\section{C3 Performance Results}

\begin{center}
\renewcommand{\arraystretch}{1.6}
\begin{tabular}{@{}L{4.0cm}ccc@{}}
    \cellcolor{dark}\textcolor{white}{\textbf{Sensor Configuration}} &
    \cellcolor{dark}\textcolor{white}{\textbf{Impl.\ Ind.}} &
    \cellcolor{dark}\textcolor{white}{\textbf{PIC (GNN)}} &
    \cellcolor{dark}\textcolor{white}{\textbf{Advantage}} \\
    Complete Visibility   & 9.10 & 3.89 & \textbf{2.34$\times$} \\
    Intersecting-Span     & 5.04 & 1.61 & \textbf{3.13$\times$} \\
    Disjoint-Span         & 5.80 & 1.84 & \textbf{3.15$\times$} \\
    Incomplete Coverage   & 4.63 & 1.05 & \textbf{4.41$\times$} \\
\end{tabular}
\end{center}

\vspace{0.4em}
\includegraphics[width=\textwidth, height=3.5cm, keepaspectratio]{eval_scale.png}

\shadedbox{Implicit Indication \textbf{consistently and substantially} outperforms
PIC across \textit{every} sensor configuration.
The advantage grows as the environment becomes \textbf{more challenging}
(incomplete coverage: $4.41\times$).}

%% ------------------------------------------------------------------
\section{C3 Robustness Comparison}\twocolumn\raggedright

\subsection{Finding}

\begin{itemize}
    \item Implicit Indication also \textbf{outperforms} PIC across the
        eight robustness evaluation conditions
    \item Both sensor degradation and sensor improvement scenarios
        favor the representational approach
    \item Zero-shot generalization to novel agent types:
        Implicit Indication handles it; PIC struggles
\end{itemize}

\subsection{Additional Baseline}

\begin{itemize}
    \item Implicit Indication also \textbf{outperforms HAPPO} under
        perturbed conditions---despite using a \textit{single shared policy}
        vs.\ HAPPO's per-type policies
    \item The simpler, more general approach wins on every front
\end{itemize}

\pagebreak

\includegraphics[width=\columnwidth, height=2.5cm, keepaspectratio]{eval_heatmaps.png}
\vspace{0.4em}
\includegraphics[width=\columnwidth, height=2.5cm, keepaspectratio]{eval_boxes.png}

%% ------------------------------------------------------------------
\section[C3 Synthesis]{C3 Synthesis: Why the Simpler Approach Wins}

\begin{itemize}
    \item PIC's GNN aggregates node features via symmetric pooling---
        appropriate when \textbf{relational structure} between agents is rich
    \item In HyperGrid, agent interactions are mediated by the task environment,
        not by direct agent-to-agent relationships
    \item Explicit \textbf{observation masking} provides a cleaner gradient signal:
        each agent's distinct channel pattern carries informative structure
        that the MLP can exploit directly
    \item GNN aggregation may \textit{smooth away} that structure
\end{itemize}

\vspace{0.5em}
\shadedbox{\textbf{Key principle:} Architectural sophistication is
\textit{not} a substitute for representational clarity.
Match the invariance mechanism to the \textbf{structure of the problem},
not to the expressive power of the architecture.}

\vspace{0.5em}
\textit{For structurally heterogeneous discrete domains with semantically
decomposable observations: \textbf{representation first, architecture second.}}


%% ==========================================================================
%% SYNTHESIS
%% ==========================================================================

\chapter{Synthesis}

%% ------------------------------------------------------------------
\section{Cross-Cutting Insights}

\begin{enumerate}
    \item \textbf{Representation Over Architecture}---Observation schemas,
        homogenized spaces, and input encodings matter more than algorithmic
        sophistication. C1, C2, and C3 all confirm this independently.

    \item \textbf{Task Structure as Critical Moderator}---No universal method exists.
        C1 shows curriculum benefit depends on coordination tightness and
        heterogeneity type. C3 shows GNN advantage depends on relational richness.

    \item \textbf{Unified Sharing Without Interchangeability}---Semantic
        decomposability enables a single policy to serve structurally
        distinct agents without assuming they are interchangeable.

    \item \textbf{Match Method to Heterogeneity Type}---Behavioral heterogeneity
        is a candidate for curriculum strategies; structural/intrinsic heterogeneity
        calls for homogenization.

    \item \textbf{Emergent Robustness}---Deployment flexibility (zero-shot
        generalization, sensor tolerance, team-size flexibility) arises
        \textit{for free} from representational design---no explicit robustness
        training required.
\end{enumerate}

%% ------------------------------------------------------------------
\section{Design Principles for HARL}

\begin{itemize}
    \item \textbf{Representation first.}
        Before choosing an algorithm or architecture, ask:
        \textit{does my observation space support the invariances I need?}

    \item \textbf{Design for semantic decomposability.}
        Observation elements should carry \textit{consistent meaning}
        across all agent types that share them.

    \item \textbf{Use masking and homogenization} for flexible team compositions;
        it is simpler and more effective than architectural alternatives
        in semantically decomposable domains.

    \item \textbf{Consider gradient interference.}
        Disjoint observation spans reduce shared-gradient noise;
        structure the observation space to minimize unintended overlap.

    \item \textbf{Choose curriculum strategies based on task structure.}
        Tight coordination and high role-stability favor pretraining;
        dynamic role heterogeneity may not benefit.
\end{itemize}

%% ------------------------------------------------------------------
\section{Contributions Summary}

\begin{center}
\renewcommand{\arraystretch}{1.6}
\begin{tabular}{@{}L{1.5cm}L{4.0cm}L{4.0cm}L{3.5cm}@{}}
    \cellcolor{dark}\textcolor{white}{\textbf{}} &
    \cellcolor{dark}\textcolor{white}{\textbf{Central Claim}} &
    \cellcolor{dark}\textcolor{white}{\textbf{Key Finding}} &
    \cellcolor{dark}\textcolor{white}{\textbf{Design Implication}} \\
    \textbf{C1} &
        Pretraining smaller teams can reduce cost &
        Benefit depends on task coordination structure and heterogeneity type &
        Obs.\ schema design determines \textit{whether} curricula can help \\
    \textbf{C2} &
        One shared policy can handle structural heterogeneity without IDs &
        Matches HAPPO; $\frac{1}{|I|}$ storage; zero-shot robust &
        Semantic decomposability is the key design requirement \\
    \textbf{C3} &
        Representational $>$ architectural invariance &
        2.3--4.4$\times$ advantage over GNN-based PIC &
        Simple masking beats complex graph aggregation in this problem class \\
\end{tabular}
\end{center}

%% ------------------------------------------------------------------
\section{The Unifying Answer}

\begin{Center}
    \large
    \textbf{DIE-HARL}\\[0.8ex]
    \normalsize
    \begin{tabular}{@{}cL{9cm}@{}}
        \textbf{\textcolor{light}{D}urable} &
            (C3) --- Representational invariance yields robustness
            to deployment variation without retraining \\[0.6ex]
        \textbf{\textcolor{light}{I}mplicit} &
            (C2) --- Observation structure encodes agent identity
            without explicit type indicators \\[0.6ex]
        \textbf{\textcolor{light}{E}xtensible} &
            (C1) --- Fixed-width schemas enable policy reuse
            across team sizes via curriculum \\
    \end{tabular}

    \vspace{1.5ex}
    \shadedbox{Training efficiency, heterogeneity handling, and deployment
    flexibility in HARL are fundamentally \textbf{representational} problems
    before they are architectural ones.}
\end{Center}


%% ==========================================================================
%% LIMITATIONS & FUTURE WORK
%% ==========================================================================

\chapter{Limitations \& Future Work}

%% ------------------------------------------------------------------
\section{Limitations}

\begin{itemize}
    \item \textbf{Single algorithm family (PPO-based)}---results may not
        generalize to off-policy methods (SAC, TD3) or evolutionary approaches

    \item \textbf{Controlled simulation environments}---real-world complexity
        (noisy sensors, communication delay, physical dynamics) not yet tested

    \item \textbf{Limited team sizes} (max 8 agents)---scaling properties
        for much larger teams remain under-explored

    \item \textbf{Semantic decomposability assumption}---not all operational
        sensor suites will satisfy this; mixed-modality observations
        (e.g., vision + LiDAR with different spatial encodings) require
        additional engineering

    \item \textbf{Predominantly discrete action spaces}---C2 and C3 evaluated
        in HyperGrid; continuous-action heterogeneity not covered
\end{itemize}

%% ------------------------------------------------------------------
\section{Future Directions}

\begin{itemize}
    \item \textbf{Action-space heterogeneity}: extending Implicit Indication
        to agents with structurally different \textit{action spaces}, not just
        observation spaces

    \item \textbf{Dynamic team composition}: agents joining or leaving the
        team mid-episode; evaluating how well homogenized policies adapt

    \item \textbf{Hybrid approaches}: combining representational and architectural
        invariance for domains with rich relational structure \textit{and}
        semantic observation structure

    \item \textbf{Real-world validation}: physical platforms with noisy sensors,
        communication constraints, and partial failures

    \item \textbf{Formal convergence guarantees}: theoretical analysis of
        convergence rates for implicit indication under CTDE training

    \item \textbf{Towards the Replicator vision}: scaling these methods to
        operationally relevant heterogeneous swarms in contested environments
\end{itemize}


%% ==========================================================================
%% CONCLUSION
%% ==========================================================================

\chapter{Conclusion}

%% ------------------------------------------------------------------
\section{Conclusion}

\begin{itemize}
    \item \textbf{C1} demonstrated that curriculum-based team scaling is
        \textit{possible and beneficial} when the observation schema supports
        policy transfer, and that task coordination structure determines
        the nature of that benefit

    \item \textbf{C2} introduced Implicit Indication: a representational
        framework enabling a single shared policy to operate over
        structurally heterogeneous agents without explicit type identifiers---
        achieving HAPPO-level performance with $\frac{1}{|I|}$ storage and
        zero-shot deployment robustness

    \item \textbf{C3} provided empirical evidence that representational
        invariance substantially outperforms architectural invariance
        (GNN-based PIC) in structurally heterogeneous discrete domains,
        with consistent advantages of 2.3--4.4$\times$

    \item Together, these contributions support a unified design philosophy:
        \textbf{invest in representation design before adding architectural
        complexity}
\end{itemize}

\vspace{0.5em}
\shadedbox{The path toward robust, efficient, and deployable heterogeneous
multi-agent systems runs through \textbf{how we design the spaces our
agents inhabit}---not just the algorithms that inhabit them.}

\closing

%% ==========================================================================
%% BACKUP SLIDES
%% ==========================================================================

\chapter[Backup]{Backup Slides}

%% ------------------------------------------------------------------
\section[B1: Iteration Cost]{B1 --- C1 Iteration Cost Analysis}\twocolumn\raggedright

\subsection{Normalized Training Cost}

\begin{itemize}
    \item How much does it actually cost to train via curriculum vs.\ from scratch?
    \item Agent-steps metric enables fair comparison across team sizes
    \item Pretraining at $k$ agents then fine-tuning to $N$ agents is cost-effective
        when $T_\text{ft} \ll T_\text{scratch}$
\end{itemize}

\pagebreak

\includegraphics[width=\columnwidth, height=5cm, keepaspectratio]{iter_cost.png}

%% ------------------------------------------------------------------
\section[B2: Deep Sets Theory]{B2 --- Theoretical Basis for Implicit Indication}

\begin{itemize}
    \item \textbf{Deep Sets} (Zaheer et al., 2017): any permutation-invariant function
        $f: 2^\mathcal{X} \to \mathbb{R}$ can be decomposed as
        $f(\mathcal{X}) = \rho\!\left(\sum_{x \in \mathcal{X}} \phi(x)\right)$
        for suitable $\phi, \rho$

    \item \textbf{Key result for C2}: for any collection of functions
        $\{f_i : \mathcal{O}_i \to \mathbb{R}\}_{i \in I}$ over
        heterogeneous input spaces, there exists a single function
        $F : \tilde{\mathcal{O}} \to \mathbb{R}$ such that
        $F(\tilde{o}) = f_i(o_i)$ for all $i \in I$

    \item \textbf{Condition}: the embedding $\phi_i : \mathcal{O}_i \to \tilde{\mathcal{O}}$
        must be \textit{injective} for each $i$, and the images
        $\phi_i(\mathcal{O}_i)$ for distinct elements must not collide
        (satisfied by disjoint-semantic channel assignment)

    \item This formally guarantees that a single neural network over
        $\tilde{\mathcal{O}}$ is \textbf{at least as expressive} as the
        collection of per-type networks
\end{itemize}

%% ------------------------------------------------------------------
\section[B3: HAPPO Detail]{B3 --- HAPPO vs.\ Implicit Indication (Detail)}\twocolumn\raggedright

\subsection{HAPPO (Heterogeneous-Agent PPO)}

\begin{itemize}
    \item Maintains \textbf{separate actor and critic} for each agent type $i \in I$
    \item Sequential update order: agents updated one type at a time,
        treating others as part of the environment
    \item Strong performance; the principled heterogeneous baseline
\end{itemize}

\subsection{The Trade-off}

\begin{itemize}
    \item Storage: $|I|$ actors + $|I|$ critics
    \item No generalization across agent types (novel type $\Rightarrow$ retrain)
    \item No zero-shot transfer to unseen team compositions
\end{itemize}

\pagebreak

\begin{center}
\renewcommand{\arraystretch}{1.5}
\begin{tabular}{@{}L{2.5cm}cc@{}}
    \cellcolor{dark}\textcolor{white}{\textbf{Property}} &
    \cellcolor{dark}\textcolor{white}{\textbf{HAPPO}} &
    \cellcolor{dark}\textcolor{white}{\textbf{Impl.\ Ind.}} \\
    Training perf. & Baseline & Matches \\
    Storage & $|I|\times$ & $1\times$ \\
    Novel type & Retrain & Zero-shot \\
    Sensor change & Retrain & Tolerant \\
    Team-size flex. & Fixed & Flexible \\
\end{tabular}
\end{center}

%% ------------------------------------------------------------------
\section[B4: HyperGrid Detail]{B4 --- HyperGrid Environment Details}\twocolumn\raggedright

\subsection{Design Goals}

\begin{itemize}
    \item Isolate the effect of \textbf{observation structure} as the
        independent variable
    \item Allow precise control over which agents share which channels
    \item Support systematic evaluation across heterogeneity configurations
\end{itemize}

\subsection{Mechanics}

\begin{itemize}
    \item $n$-dimensional discrete grid; agents collect items
    \item Each cell has multiple \textbf{channel layers}
        (analogous to sensor modalities)
    \item Agent type determines which channels are populated in its observation
    \item Task reward is cooperative; agents must coordinate based on
        complementary observations
\end{itemize}

\pagebreak

\includegraphics[width=\columnwidth, height=4.5cm, keepaspectratio]{marl-minigrid.png}

\scriptsize\textit{Color-coded channels in the HyperGrid environment.
Agent type determines which channels are visible.}

%% ------------------------------------------------------------------
\section[B5: C2 Eval Detail]{B5 --- C2 Robustness Evaluation (Detail)}\twocolumn\raggedright

\subsection{Eight Conditions Explained}

\begin{enumerate}
    \item \textbf{Sensor loss}: remove a channel from agents that had it
    \item \textbf{Sensor degradation}: add Gaussian noise to a channel
    \item \textbf{Sensor improvement}: remove noise from a previously noisy channel
    \item \textbf{Coverage reduction}: reduce total number of populated channels
    \item \textbf{Coverage expansion}: add a new channel to some agents
    \item \textbf{Shuffled assignments}: swap which agents get which channels
    \item \textbf{Incomplete generalization}: test on incomplete coverage config
    \item \textbf{Zero-shot}: introduce an agent type never seen in training
\end{enumerate}

\pagebreak

\includegraphics[width=\columnwidth, height=5cm, keepaspectratio]{eval_scale.png}

%% ------------------------------------------------------------------
\section[B6: Disjoint Span]{B6 --- Why Disjoint-Span Training is Strongest}

\begin{itemize}
    \item \textbf{Disjoint-span}: no agent type shares channels with any other
    \item This creates \textbf{maximally informative observation patterns}:
        the populated/masked structure uniquely identifies the agent type
    \item In intersecting-span configurations, some channels are shared---
        the gradient signal for those elements is averaged across agent types,
        potentially introducing \textit{gradient interference}
    \item Disjoint spans \textbf{eliminate gradient interference}:
        each channel's gradient comes exclusively from agents of one type
    \item \textbf{Design implication}: where possible, prefer observation
        designs with clean separation between agent-specific and shared elements
\end{itemize}

\vspace{0.5em}
\shadedbox{The most semantically decomposable observation structure
also yields the cleanest learning signal.
\textbf{Good representational design has cascading benefits.}}

%% ------------------------------------------------------------------
\section[B7: Smit Comparison]{B7 --- Comparison to Prior Curriculum Work}

\begin{itemize}
    \item Smit et al.\ (2023) investigate curriculum-based team scaling in
        football environments---the most directly comparable prior work
    \item Their environments feature \textbf{fixed spatial roles} and
        strong coordination requirements
    \item Their findings: curriculum improves performance in high-coordination
        settings, consistent with C1's Multiwalker results
    \item \textbf{C1 extends} this finding by:
    \begin{itemize}
        \item Evaluating across multiple environment types
        \item Introducing the agent-steps normalized cost metric
        \item Explicitly identifying the observation-space coupling problem
            as a prerequisite concern
        \item Testing on LBF, which exposes the limits of the curriculum approach
    \end{itemize}
\end{itemize}

\vspace{0.5em}
\includegraphics[width=0.5\textwidth, height=3cm, keepaspectratio]{smit_football.png}

%% ------------------------------------------------------------------
\section[B8: Action-Space Heterogeneity]{B8 --- Future: Action-Space Heterogeneity}

\begin{itemize}
    \item C2 and C3 address \textbf{observation-space} structural heterogeneity:
        agents differ in what they \textit{see}
    \item The natural next extension: \textbf{action-space heterogeneity}---
        agents differ in what they \textit{can do}
    \item Example: a UAV can fly; a ground vehicle can drive; both must coordinate
    \item Challenges:
    \begin{itemize}
        \item A homogenized action space is non-trivial---invalid actions
            must be masked or penalized
        \item The policy must learn that certain actions are
            \textit{structurally unavailable} to it, not just suboptimal
        \item Multi-discrete and hybrid continuous-discrete action spaces
            add further complexity
    \end{itemize}
    \item Preliminary hypothesis: the same semantic decomposability principle
        should extend---but the masking must occur at both \textit{input} and
        \textit{output} stages
\end{itemize}

\end{document}
